% ==================================================================
% id:       lim-037
% topic:    גבולות ורציפות של פונקציות
% subtopic: אריתמטיקה של גבולות
% solution: full
% ==================================================================
\begin{enumerate}
\item מצאו את כל הערכים של $a,b$ עבורם $\ds\lim_{x\to2}\frac{x^2+ax+b}{x-2}=5$.
\ifsol
\hint{אם הגבול קיים, לאן חייב לשאוף המונה? כתבו אותו כמכפלה של המנה ב-$x-2$.}
\begin{solbox}
נניח שהגבול שווה ל-$5$. לכל $x\ne2$
\[x^2+ax+b=\frac{x^2+ax+b}{x-2}\cdot(x-2)\]
ולכן
\[\lim_{x\to2}\left(x^2+ax+b\right)=\lim_{x\to2}\left(\frac{x^2+ax+b}{x-2}\cdot(x-2)\right)=5\cdot0=0\]
מצד שני, המונה הוא פולינום, ולכן הגבול שלו ב-$2$ הוא $4+2a+b$. מיחידות הגבול $4+2a+b=0$, כלומר $b=-4-2a$. אז
\[x^2+ax-4-2a=(x-2)(x+2)+a(x-2)=(x-2)(x+2+a)\]
ולכן לכל $x\ne2$ המנה שווה ל-$x+2+a$, ולכן הגבול שלה ב-$2$ הוא $4+a$. מכאן $4+a=5$, כלומר
\[a=1,\qquad b=-6\]
\end{solbox}
\fi
\item מצאו את כל הערכים של $a$ עבורם הגבול $\ds\lim_{x\to1}\frac{\sqrt{x+a}-2}{x-1}$ קיים (וסופי), וחשבו אותו.
\ifsol
\hint{כמו בסעיף הקודם, המונה חייב לשאוף ל-$0$. אחר כך כפלו בצמוד.}
\begin{solbox}
כדי שהביטוי יהיה מוגדר בסביבה נקובה של $1$ צריך $1+a\ge0$. אם הגבול $L$ קיים, אז כמו בסעיף הקודם
\[\sqrt{x+a}-2=\frac{\sqrt{x+a}-2}{x-1}\cdot(x-1)\]
ולכן
\[\lim_{x\to1}\left(\sqrt{x+a}-2\right)=\lim_{x\to1}\left(\frac{\sqrt{x+a}-2}{x-1}\cdot(x-1)\right)=L\cdot0=0\]
ומצד שני הגבול של המונה ב-$1$ הוא $\sqrt{1+a}-2$. לכן $\sqrt{1+a}=2$, כלומר $a=3$. עבור $a=3$ נכפול בצמוד:
\[\frac{\sqrt{x+3}-2}{x-1}=\frac{x-1}{(x-1)\lr{\sqrt{x+3}+2}}=\frac{1}{\sqrt{x+3}+2}\]
ולכן
\[\lim_{x\to1}\frac{\sqrt{x+3}-2}{x-1}=\lim_{x\to1}\frac{1}{\sqrt{x+3}+2}=\frac14\]
\end{solbox}
\fi
\item יהי $n$ טבעי. חשבו את $\ds\lim_{x\to1}\frac{x^n-1}{x-1}$.
\ifsol
\hint{$x^n-1=(x-1)\lr{x^{n-1}+x^{n-2}+\dots+x+1}$.}
\begin{solbox}
לכל $x\ne1$
\[\frac{x^n-1}{x-1}=x^{n-1}+x^{n-2}+\dots+x+1\]
זה פולינום, ולכן הגבול שלו ב-$1$ הוא הערך שלו ב-$1$: סכום של $n$ מחוברים ששווים ל-$1$. הגבול הוא $n$.
\end{solbox}
\fi
\end{enumerate}
